\section{Discusión}

\subsection{Organización modular de la red}

Los genes asociados al término HP:0002647 no se distribuyeron al azar en la red de interacción, sino que se agruparon en seis módulos, dos de ellos de mayor tamaño. Esta organización concuerda con el concepto de módulo de enfermedad, según el cual las proteínas implicadas en un mismo fenotipo tienden a interaccionar entre sí y a ocupar una vecindad concreta de la red \cite{Barabasi2011}. Además, los módulos reprodujeron los tres procesos que la literatura relaciona con la enfermedad aórtica torácica hereditaria: la matriz extracelular, la señalización de TGF-$\beta$ y la contracción del músculo liso \cite{Pinard2019,Lindsay2011}. Este resultado es, por tanto, el esperado y sirve sobre todo para confirmar que la red recoge la biología conocida de la enfermedad. Su interés reside más bien en cómo se relacionan estos procesos entre sí y en los módulos que no encajan en este esquema.

Los dos módulos principales, de diez genes cada uno, correspondieron a la matriz extracelular y las fibras elásticas (\textit{BGN}, \textit{COL3A1}, \textit{EFEMP2}, \textit{ELN}, \textit{EMILIN1}, \textit{FBN1}, \textit{LOX}, \textit{MFAP5}, \textit{P4HA2} y \textit{PLOD1}) y a la señalización de TGF-$\beta$/SMAD (\textit{ENG}, \textit{SMAD2}, \textit{SMAD3}, \textit{SMAD4}, \textit{SMAD6}, \textit{TGFB2}, \textit{TGFB3}, \textit{TGFBR1}, \textit{TGFBR2} y \textit{TGFBR3}). Aunque el algoritmo de Louvain los separó, ambos procesos están estrechamente ligados: los modelos animales de las formas hereditarias mostraron la implicación inesperada de la vía de TGF-$\beta$ en enfermedades causadas por defectos de proteínas estructurales de la matriz, como la fibrilina-1 en el síndrome de Marfan \cite{Lindsay2011}. En consonancia con ello, los tres genes con más interacciones de la red, \textit{TGFBR1}, \textit{TGFBR2} y \textit{FBN1}, pertenecieron a estos dos módulos, lo que sitúa a la matriz extracelular y a la señalización de TGF-$\beta$ en el núcleo de la red.

El módulo de músculo liso agrupó a \textit{ACTA2}, \textit{MYH11} y \textit{MYLK}, que codifican componentes del aparato contráctil y de su regulación. Estas proteínas forman parte de la unidad elastina-contráctil, la estructura que une las fibras elásticas con las células de músculo liso de la capa media y cuya alteración es común a muchos de los genes que causan aneurismas y disecciones de la aorta torácica \cite{Karimi2016}. Llama la atención que \textit{PRKG1}, que codifica una quinasa de la misma vía y cuenta con evidencia fuerte de asociación con la enfermedad \cite{Pinard2019,Renard2018} quedara sin interacciones con el resto de genes con el umbral utilizado. Este caso muestra que la posición de un gen en la red no refleja necesariamente su relevancia clínica.

El módulo de Notch (\textit{NOTCH1}, \textit{HEY2} y \textit{NKX2-5}) apunta a un mecanismo distinto, relacionado con el desarrollo cardiovascular más que con el mantenimiento de la pared aórtica. Las mutaciones en \textit{NOTCH1} causan un espectro de anomalías del desarrollo de la válvula aórtica, y los represores transcripcionales de la familia de \textit{HEY2}, activados por la señalización de Notch1, intervienen en ese mismo proceso \cite{Garg2005}.

\subsection{Módulos inesperados}

El módulo formado por \textit{HLA-B}, \textit{HLA-DRB1} y \textit{PTPN22} fue el resultado más inesperado, ya que reúne genes del sistema inmunitario sin relación con la estructura de la pared aórtica. Según las anotaciones de HPO, los tres genes se asociaron al término a través de una misma enfermedad, la arteritis de células gigantes (ORPHA:397), una vasculitis de grandes vasos. En esta enfermedad, la región HLA concentra la mayor parte de la susceptibilidad genética, con contribuciones de \textit{HLA-DRB1} y de \textit{HLA-B} \cite{Carmona2015}, y la variante R620W de \textit{PTPN22} constituye otro factor de riesgo \cite{Serrano2013}. Que este módulo no se conectara con los de matriz extracelular y TGF-$\beta$ es coherente con una vía de daño de la pared aórtica distinta, de origen inflamatorio y no estructural.

Del mismo modo, \textit{ABCC6} y \textit{ENPP1}, ambos asociados al término a través de la calcificación arterial generalizada de la infancia (ORPHA:51608), formaron un módulo aislado del resto. Las mutaciones en cualquiera de los dos genes pueden causar la calcificación arterial generalizada de la infancia y el pseudoxantoma elástico, dos extremos de un mismo espectro de calcificación ectópica que afecta a la lámina elástica interna de las arterias \cite{Nitschke2012}. Este módulo representaría, por tanto, una tercera vía de daño de las fibras elásticas, a través de su calcificación.

\subsection{Enriquecimiento funcional y elección del fondo}

El análisis de enriquecimiento de la lista completa, con el fondo de los 5277 genes con anotación en HPO, identificó 410 procesos biológicos de GO, encabezados por \textit{extracellular matrix assembly}, y el del módulo de TGF-$\beta$, 81 procesos, entre ellos \textit{SMAD protein signal transduction}. Ambos resultados confirmaron la interpretación biológica de los módulos.

En cambio, el módulo de matriz extracelular no mostró ningún término enriquecido. Este resultado no indica que el módulo carezca de una función común, sino que se debe al fondo utilizado: en el análisis de cada módulo, el fondo fueron los 31 genes de la red, que ya son en buena parte proteínas de la matriz extracelular. Como el test de enriquecimiento mide si una función está más representada en la lista que en el fondo, un módulo cuya función coincide con la del fondo no puede destacar. Esto ilustra que la elección del fondo condiciona el resultado del enriquecimiento, una fuente de sesgo descrita en la literatura que a menudo se pasa por alto \cite{Timmons2015}. Por el mismo motivo, para la lista completa se utilizó como fondo el conjunto de genes con alguna anotación en HPO y no todos los genes con anotación en GO: como la lista procede de HPO, compararla con todo el genoma habría inflado el número de términos significativos. Un número menor de términos con el fondo adecuado es, por tanto, un resultado más fiable.

Por último, solo se enriquecieron los módulos de cinco genes o más, ya que en listas tan pequeñas unos pocos genes bastan para que un término resulte significativo. Por ello, los módulos pequeños (músculo liso, Notch, HLA/\textit{PTPN22} y \textit{ABCC6}/\textit{ENPP1}) se interpretaron únicamente a partir de la literatura.